% =============================================================================
%                    WEEK 7: FAIR DIVISION (PART I)
% =============================================================================
\section{Week 7: Fair Division (Part I)}

% -----------------------------------------------------------------------------
%                              7.1 TOPIC SUMMARY
% -----------------------------------------------------------------------------
\subsection{Topic Summary}

\begin{keybox}[Learning Objectives]
By the end of this week, you will be able to:
\begin{enumerate}
  \item Explain foundational concepts of fair division: \textbf{proportionality}, \textbf{envy-freeness}, and \textbf{efficiency}.
  \item Model resource allocation problems as fair division scenarios using mathematical and algorithmic frameworks.
  \item Apply classic algorithms (\textbf{Cut-and-Choose}, \textbf{Dubins--Spanier}, \textbf{Even--Paz}, \textbf{Selfridge--Conway}) to achieve fair outcomes for \textbf{divisible goods}.
  \item Analyse the strengths and limitations of different fair division procedures (guarantees, complexity, strategic aspects).
  \item Assess allocation solutions for fairness and optimality using formal properties and counterexamples.
  \item Solve practical and theoretical problems involving cake-cutting, rent division, and item allocation.
\end{enumerate}
\end{keybox}

\separator

\subsubsection*{The Big Picture}

The first half of the module covered strategic interaction (non-cooperative and cooperative games). The second half turns to \textbf{collective decision-making}: fair division (Weeks~7--8), mechanism design, and voting. Fair division asks: how should one or more goods be divided among agents so that every agent considers the outcome \emph{fair}? This week focuses on \textbf{divisible goods} (cake-cutting); Week~8 treats \textbf{indivisible goods}.

\separator

\subsubsection*{What Examiners Like to Ask}

\begin{itemize}
  \item Define proportionality and envy-freeness; prove one implies the other for $n=2$.
  \item Trace the \textbf{Cut-and-Choose}, \textbf{Dubins--Spanier}, or \textbf{Even--Paz} algorithm on a concrete instance.
  \item Show that a given algorithm is (or is not) proportional / envy-free / equitable.
  \item Analyse variants of moving-knife procedures (e.g.\ ``$1/3$-envy-free'' bounds).
  \item Compute the \textbf{price of proportionality} (welfare loss from requiring proportionality).
  \item Determine whether a connected allocation with given value guarantees exists.
  \item State complexity results under the \textbf{Robertson--Webb model}.
\end{itemize}

\separator

% -----------------------------------------------------------------------------
%           7.2 OVERVIEW + CORE CONCEPTS + EXTENDED THEORY
% -----------------------------------------------------------------------------
\subsection{Overview + Core Concepts + Extended Theory}

\subsubsection*{Roadmap}
\begin{enumerate}
  \item Setting: agents, goods, divisible vs.\ indivisible
  \item Fairness criteria: proportionality, envy-freeness, equitability, Pareto-optimality
  \item Cake-cutting model and the Robertson--Webb computational model
  \item Algorithm: Cut-and-Choose (2 agents, EF)
  \item Algorithm: Dubins--Spanier moving knife ($n$ agents, proportional)
  \item Algorithm: Even--Paz ($n$ agents, proportional, optimal complexity)
  \item Algorithm: Selfridge--Conway (3 agents, EF)
  \item General existence and complexity of envy-free algorithms
  \item Worked examples and exam-style problems
\end{enumerate}

\bigseparator

% --- Setting ---
\subsubsection{The Fair Division Setting}

\begin{defbox}[Fair Division Problem]
A \textbf{fair division problem} consists of:
\begin{itemize}
  \item A finite set of \textbf{agents} $N = \{1, \ldots, n\}$.
  \item A set of \textbf{goods} (resources/items) to be divided.
  \item Each agent $i$ has a \textbf{valuation function} $v_i$ over (portions of) the goods.
\end{itemize}
An \textbf{allocation} assigns (parts of) the goods to agents.
\end{defbox}

\textbf{Types of goods:}
\begin{itemize}
  \item \textbf{Divisible} (e.g.\ cake): any fraction can be allocated. This week's focus.
  \item \textbf{Indivisible} (e.g.\ a book): must be allocated whole. Week~8's focus.
  \item \textbf{Mixed}: indivisible goods + money (a divisible resource).
\end{itemize}

We assume goods are \textbf{static} (properties don't change), available in \textbf{single units}, and \textbf{cannot be shared} between agents.

\separator

% --- Fairness Criteria ---
\subsubsection{Fairness Criteria}

Let $(A_1, \ldots, A_n)$ be an allocation to $n$ agents.

\begin{defbox}[Proportionality]
An allocation is \textbf{proportional} if every agent receives at least their ``fair share'':
\[
  v_i(A_i) \;\ge\; \frac{1}{n} \quad \text{for all } i \in N.
\]
(Valuations are normalised so $v_i(\text{whole cake}) = 1$.)
\end{defbox}

\begin{defbox}[Envy-Freeness (EF)]
An allocation is \textbf{envy-free} if no agent prefers another agent's bundle:
\[
  v_i(A_i) \;\ge\; v_i(A_j) \quad \text{for all } i, j \in N.
\]
\end{defbox}

\begin{thmbox}[EF $\Rightarrow$ Proportionality]
For $n$ agents, envy-freeness implies proportionality.

\textit{Proof.} If the allocation is EF, then $v_i(A_i) \ge v_i(A_j)$ for all $j$. Summing over all $j$:
$n \cdot v_i(A_i) \ge \sum_{j=1}^{n} v_i(A_j) = v_i(\text{whole cake}) = 1$,
so $v_i(A_i) \ge 1/n$. \qed
\end{thmbox}

\begin{thmbox}[Equivalence for $n=2$]
For two agents, an allocation is envy-free \textbf{if and only if} it is proportional.

\textit{Proof.} EF $\Rightarrow$ proportional: shown above. Proportional $\Rightarrow$ EF: if $v_i(A_i) \ge 1/2$, then $v_i(A_j) = 1 - v_i(A_i) \le 1/2 \le v_i(A_i)$. \qed
\end{thmbox}

\begin{tipbox}[Exam Tip]
For $n \ge 3$, proportionality does \textbf{not} imply envy-freeness. Be ready to give a counterexample (e.g.\ Dubins--Spanier chocolate/sponge example below).
\end{tipbox}

\begin{defbox}[Equitability]
An allocation is \textbf{equitable} if all agents derive the same value: $v_i(A_i) = v_j(A_j)$ for all $i, j$.
\end{defbox}

\begin{defbox}[Pareto Optimality]
An allocation is \textbf{Pareto optimal} if there is no other allocation that makes some agent strictly better off without making any agent worse off.
\end{defbox}

\separator

% --- Cake Cutting Model ---
\subsubsection{Cake-Cutting Model}

The ``cake'' is the interval $[0, 1]$. Each agent $i$ has a valuation $v_i$ that is:
\begin{itemize}
  \item \textbf{Non-negative}: $v_i(X) \ge 0$ for all $X \subseteq [0,1]$.
  \item \textbf{Additive}: $v_i(X \cup Y) = v_i(X) + v_i(Y)$ for disjoint $X, Y$.
  \item \textbf{Divisible}: for any $X$ with $v_i(X) > 0$ and $0 < \lambda < 1$, there exists $Y \subseteq X$ with $v_i(Y) = \lambda \cdot v_i(X)$.
  \item \textbf{Normalised}: $v_i([0,1]) = 1$.
\end{itemize}

\begin{defbox}[Robertson--Webb Model]
The complexity of cake-cutting algorithms is measured by the number of the following two queries:
\begin{enumerate}
  \item \textbf{Evaluation query} $\mathrm{Eval}_i(x, y)$: returns $v_i([x, y])$ --- ``How much does agent $i$ value piece $[x,y]$?''
  \item \textbf{Cut query} $\mathrm{Cut}_i(x, \alpha)$: returns $y$ such that $v_i([x, y]) = \alpha$ --- ``Starting at $x$, where must we cut so agent $i$ values the piece at exactly $\alpha$?''
\end{enumerate}
\end{defbox}

\separator

% --- Cut and Choose ---
\subsubsection{Cut-and-Choose (2 Agents)}

\begin{thmbox}[Cut-and-Choose Algorithm]
\begin{enumerate}
  \item Player~1 cuts the cake into two pieces $X_1, X_2$ with $v_1(X_1) = v_1(X_2) = 1/2$.
  \item Player~2 chooses the piece they prefer; Player~1 gets the other.
\end{enumerate}
\textbf{Properties:}
\begin{itemize}
  \item \textbf{Envy-free} (and hence proportional for $n=2$): Player~1 values both pieces equally, so never envies. Player~2 picks their favourite, so never envies.
  \item \textbf{Not necessarily equitable.}
  \item \textbf{Complexity}: 2 operations (1 cut + 1 eval) under Robertson--Webb.
\end{itemize}
\end{thmbox}

\begin{exbox}[Inequity in Cut-and-Choose]
Cake: left half chocolate, right half sponge. Agent~1 values only chocolate ($v_1 = 1$ for chocolate half). Agent~2 values everything uniformly.

\textbf{Agent~1 cuts first:} Cuts chocolate half in two. Agent~2 picks one quarter (value $1/4$ to them) + the sponge side is not offered this way. Actually: Agent~1 must cut so \emph{they} value both halves equally. Since Agent~1 values only chocolate, they cut at the midpoint of the chocolate region, i.e.\ at $x = 1/4$. Pieces: $[0, 1/4]$ and $[1/4, 1]$.
\begin{itemize}
  \item Agent~2 picks $[1/4, 1]$ (value $3/4$). Agent~1 gets $[0, 1/4]$ (value $1/2$).
\end{itemize}

\textbf{Agent~2 cuts first:} Cuts at $x = 1/2$ (values both halves equally). Agent~1 picks $[0, 1/2]$ (value $1$ --- all the chocolate!). Agent~2 gets $[1/2, 1]$ (value $1/2$).

Both are envy-free, but the values differ: the cutter advantage varies by who cuts.
\end{exbox}

\separator

% --- Dubins-Spanier ---
\subsubsection{Dubins--Spanier Moving Knife ($n$ Agents, Proportional)}

\begin{thmbox}[Dubins--Spanier Algorithm]
A referee moves a knife continuously from left to right across the cake.
\begin{enumerate}
  \item When any player finds the piece to the left of the knife worth $\ge 1/n$ to them, they shout ``Stop!''
  \item That player receives the piece to the left, leaves the game. (Ties broken arbitrarily.)
  \item Repeat with remaining $n-1$ players on the remaining cake.
  \item The last player receives the remainder.
\end{enumerate}
\end{thmbox}

\textbf{Proportionality proof:}
\begin{itemize}
  \item Any player who shouts ``Stop!'' receives a piece worth $\ge 1/n$ to them.
  \item The last player never shouted (or tied): they valued each of the $n-1$ removed pieces at $\le 1/n$, so the remainder is worth $\ge 1 - (n-1)/n = 1/n$.
\end{itemize}

\textbf{Not envy-free for $n \ge 3$:}

\begin{exbox}[Dubins--Spanier Fails Envy-Freeness]
Cake: $[0, 1/2]$ chocolate, $[1/2, 1]$ sponge. Three agents:
\begin{itemize}
  \item Agent~1: values everything uniformly.
  \item Agents~2,~3: value only sponge.
\end{itemize}
Agent~1 shouts at $1/3$, gets $[0, 1/3]$. Agents~2 and 3 both shout at $2/3$ (the point where sponge value reaches $1/3$). Say Agent~2 wins the tie, gets $[1/3, 2/3]$. Agent~3 gets $[2/3, 1]$.

Agent~2 \textbf{envies} Agent~3: Agent~3 has more sponge ($[2/3, 1]$ is all sponge, while $[1/3, 2/3]$ is half chocolate, half sponge).
\end{exbox}

\textbf{Complexity:} $\Theta(n^2)$ under Robertson--Webb (at each round, $O(n)$ queries for $n$ rounds total: $n + (n-1) + \cdots + 2 = \Theta(n^2)$).

\textbf{Also not necessarily equitable or Pareto-optimal.}

\separator

% --- Even-Paz ---
\subsubsection{Even--Paz Algorithm ($n$ Agents, Proportional, Optimal)}

\begin{thmbox}[Even--Paz Algorithm]
Divide-and-conquer approach (described for $n = 2^k$; extends to arbitrary $n$):
\begin{enumerate}
  \item Let $[s, t]$ be the current piece. If $n = 1$: give it to the remaining agent.
  \item Each agent $i$ marks a point $z_i$ on $[s, t]$ such that $v_i([s, z_i]) = \tfrac{1}{2} v_i([s, t])$.
  \item Find the median mark $z^*$ that splits the agents into two halves.
  \item Recursively allocate $[s, z^*]$ to the left half of agents and $[z^*, t]$ to the right half.
\end{enumerate}
\end{thmbox}

\textbf{Proportionality proof (by induction):}
\begin{itemize}
  \item \textbf{Base case} ($k=0$): one agent gets the whole cake, valued at $1 \ge 1/n$.
  \item \textbf{Inductive step}: at step $k$, agents sharing a piece value it at $\ge 1/2^k$. At step $k+1$, each agent is assigned to the half they value at $\ge \tfrac{1}{2} \cdot \tfrac{1}{2^k} = \tfrac{1}{2^{k+1}}$.
  \item After $\log n$ steps: each agent's piece is valued at $\ge 1/2^{\log n} = 1/n$.
\end{itemize}

\textbf{Complexity:} $O(n \log n)$ under Robertson--Webb ($\log n$ levels, $O(n)$ queries per level).

\begin{thmbox}[Optimality of Even--Paz]
Any proportional cake-cutting algorithm requires $\Omega(n \log n)$ operations under the Robertson--Webb model. Hence Even--Paz is \textbf{optimal}.
\end{thmbox}

\begin{tipbox}[Exam Tip: Comparing Proportional Algorithms]
\begin{center}
\begin{tabular}{lccc}
\toprule
\textbf{Algorithm} & \textbf{Agents} & \textbf{Guarantee} & \textbf{Complexity} \\
\midrule
Cut-and-Choose & 2 & EF (+ prop.) & $O(1)$ \\
Dubins--Spanier & $n$ & Proportional & $\Theta(n^2)$ \\
Even--Paz & $n$ & Proportional & $\Theta(n \log n)$ \\
\bottomrule
\end{tabular}
\end{center}
Dubins--Spanier is simpler but suboptimal in query complexity.
\end{tipbox}

\separator

% --- Selfridge-Conway ---
\subsubsection{Selfridge--Conway Algorithm (3 Agents, Envy-Free)}

\begin{thmbox}[Selfridge--Conway Algorithm]
Achieves \textbf{envy-free allocation for 3 agents} in a bounded number of steps.

\textbf{Stage 0 --- Initial cut and trim:}
\begin{enumerate}
  \item Player~1 cuts the cake into 3 pieces they value equally (each worth $1/3$).
  \item Player~2 examines the pieces. If the two largest (in Player~2's view) are not equal, Player~2 \textbf{trims} the largest so the two largest are tied.
  \item The trimmings become \textbf{Cake~2}; the three pieces form \textbf{Cake~1}.
\end{enumerate}

\textbf{Stage 1 --- Allocate Cake~1:}
\begin{enumerate}
  \item Player~3 picks their favourite piece from Cake~1.
  \item If Player~3 did \emph{not} take the trimmed piece $\Rightarrow$ Player~2 must take the trimmed piece.\\
        If Player~3 \emph{did} take the trimmed piece $\Rightarrow$ Player~2 picks their favourite from the remaining two.
  \item Player~1 gets the remaining piece.
\end{enumerate}
Let $T$ = the player (2 or 3) who got the trimmed piece; $T'$ = the other.

\textbf{Stage 2 --- Allocate Cake~2 (the trimmings):}
\begin{enumerate}
  \item Player $T'$ cuts Cake~2 into 3 equal pieces (in their view).
  \item Players choose in order: $T$, then Player~1, then $T'$.
\end{enumerate}
\end{thmbox}

\textbf{Envy-freeness proof (per player, using additivity to argue on Cake~1 and Cake~2 separately):}

\begin{itemize}
  \item \textbf{Player $T'$}: On Cake~2, divided it equally $\Rightarrow$ EF. On Cake~1: if $T' = 3$, they chose first among all three pieces $\Rightarrow$ EF. If $T' = 2$, they trimmed so the two largest pieces are equal, and they got one of those $\Rightarrow$ EF.

  \item \textbf{Player $T$}: On Cake~2, they choose first $\Rightarrow$ EF. On Cake~1, same argument as above.

  \item \textbf{Player~1}: On Cake~1, they cut equally so every piece is worth $1/3$ $\Rightarrow$ EF (they didn't get the trimmed piece). On Cake~2: Player~1 doesn't envy $T'$ (Player~1 chose before $T'$). Player~1 doesn't envy $T$ because $T$ got the trimmed piece on Cake~1, which is worth $\le 1/3$ to Player~1; even if $T$ gets \emph{all} of Cake~2, $T$'s total $\le 1/3$, so Player~1's ``irrevocable advantage'' ensures no envy.
\end{itemize}

\separator

% --- General EF results ---
\subsubsection{General Envy-Free Cake Cutting}

\begin{thmbox}[Existence (Brams--Taylor, 1995)]
An envy-free cake-cutting algorithm exists for \textbf{any number of players} $n$ under the Robertson--Webb model. However, its complexity cannot be bounded as a function of $n$ alone (it depends on the valuation functions).
\end{thmbox}

\begin{thmbox}[Bounded Protocol (Aziz--Mackenzie, 2016)]
A \textbf{bounded} envy-free algorithm exists for $n$ agents, but with extremely high complexity:
\[
  O\!\bigl(\underbrace{n^{n^{n^{\cdot^{\cdot^{\cdot^n}}}}}}_{5 \text{ levels}}\bigr).
\]
\end{thmbox}

\begin{thmbox}[Lower Bound (Procaccia, 2009)]
Any envy-free cake-cutting algorithm requires $\Omega(n^2)$ queries.
\end{thmbox}

\begin{tipbox}[Exam Tip: EF Complexity Summary]
\begin{center}
\begin{tabular}{lcc}
\toprule
\textbf{Agents} & \textbf{Algorithm} & \textbf{Complexity} \\
\midrule
$n = 2$ & Cut-and-Choose & $O(1)$ \\
$n = 3$ & Selfridge--Conway & $O(1)$ (bounded) \\
$n \ge 4$ & Aziz--Mackenzie & $O(n^{n^{n^{n^n}}})$ \\
\midrule
\multicolumn{2}{l}{\textbf{Lower bound (any $n$)}} & $\Omega(n^2)$ \\
\bottomrule
\end{tabular}
\end{center}
There is a large gap between the upper and lower bounds for $n \ge 4$.
\end{tipbox}

\bigseparator

% --- Worked Examples ---
\subsubsection{Worked Examples}

\begin{exbox}[Example 1: Moving Knife Variant --- $1/3$-Envy-Freeness]
\textbf{Problem (Exercise 1).} Consider a variant of the moving knife where agents shout ``Stop!'' when the left piece is worth exactly $1/3$. The piece goes to the shouting agent, and the process continues. Prove the resulting allocation is $1/3$-envy-free, i.e.\ $v_i(A_i) \ge v_i(A_j) - 1/3$ for all $i, j$.

\textbf{Solution.}
An agent who shouts ``Stop!'' receives a piece worth exactly $1/3$ to them. Any piece to the \emph{left} of theirs was cut when some other agent valued it at $1/3$; since the shouter waited, they value that earlier piece at $\le 1/3$. So no envy towards earlier agents.

For pieces to the \emph{right}: an agent might value a later piece at up to $2/3$ (the entire remaining cake was worth $2/3$ when they shouted). So the maximum envy is $v_i(A_j) - v_i(A_i) \le 2/3 - 1/3 = 1/3$. Hence $v_i(A_i) \ge v_i(A_j) - 1/3$.

For the last agent(s) who didn't shout: they valued each removed piece at $\le 1/3$, so the remainder is $\ge 1/3$, and the same bound applies.

\textbf{Is $1/3$ tight?} Yes. Consider 2 agents, cake uniform. Agent~1 shouts at $1/3$, gets $[0, 1/3]$ (value $1/3$). Agent~2 gets $[1/3, 1]$ (value $2/3$). Agent~1 envies Agent~2 by exactly $2/3 - 1/3 = 1/3$.
\end{exbox}

\separator

\begin{exbox}[Example 2: Connected Allocations]
\textbf{Problem (Exercise 2).} Two agents, Alice and Bob.
\begin{enumerate}
  \item[(a)] Is there always a connected allocation with $v_A(A_A) \ge 0.6$ and $v_B(A_B) \ge 0.4$?
  \item[(b)] What if we drop connectivity?
  \item[(c)] Is there always a connected allocation where \emph{one} agent gets $\ge 0.6$ and the other $\ge 0.4$?
\end{enumerate}

\textbf{Solution.}

\textbf{(a) No.} Counterexample: both agents value the cake uniformly. Any connected allocation is a split at some point $x$: one gets $[0,x]$, the other $[x,1]$. For Alice to get $\ge 0.6$, she needs a piece of length $\ge 0.6$, leaving Bob $\le 0.4$. But requiring $v_B \ge 0.4$ means Bob gets exactly $0.4$ only if $x = 0.6$ or $x = 0.4$. This works for uniform valuations. However, consider: Alice values $[0, 0.5]$ at $1$ and $[0.5, 1]$ at $0$; Bob values $[0, 0.5]$ at $1$ and $[0.5, 1]$ at $0$. A connected split at $x$: Alice gets $[0, x]$, Bob gets $[x, 1]$. For Alice $\ge 0.6$: need $x \ge 0.3$ (since $v_A([0,x]) = \min(2x, 1)$ for $x \le 0.5$; so $2x \ge 0.6 \Rightarrow x \ge 0.3$). For Bob $\ge 0.4$: need $v_B([x,1]) \ge 0.4$, i.e.\ $v_B([x, 0.5]) = 2(0.5 - x) \ge 0.4 \Rightarrow x \le 0.3$. So $x = 0.3$ works: Alice gets $0.6$, Bob gets $0.4$. This particular example works, but we can also swap who gets which side.

Actually, for the specific question: with \emph{arbitrary} valuations and a \emph{fixed assignment} (Alice gets left, Bob gets right), it is \textbf{not always} possible. Consider Alice values $[0.5, 1]$ at $1$ and Bob values $[0, 0.5]$ at $1$. Any connected split $[0, x]$ to Alice, $[x, 1]$ to Bob: Alice needs $v_A([0,x]) \ge 0.6 \Rightarrow x > 0.5$, so Bob gets $[x, 1]$ with $x > 0.5$, but $v_B([x,1]) = 0 < 0.4$.

\textbf{(b) Yes} (without connectivity). By a probabilistic / measure-theoretic argument, for any two additive valuations summing to $1$, we can find a (possibly disconnected) allocation giving each agent $\ge 1/2$, which exceeds both $0.6$ and $0.4$... Actually, we need Alice $\ge 0.6$ \emph{and} Bob $\ge 0.4$. Since EF gives each $\ge 1/2 > 0.4$, but Alice needs $0.6$. Without connectivity, we can partition $[0,1]$ into intervals and assign optimally. The answer is \textbf{yes}: we can always achieve $v_A \ge 0.6$ and $v_B \ge 0.4$ without connectivity (this follows from Lyapunov's theorem on the range of vector measures).

\textbf{(c) Yes.} We can always find a connected allocation where one agent gets $\ge 0.6$ and the other gets $\ge 0.4$. Consider the function $f(x) = v_A([0, x])$: it is continuous, $f(0)=0$, $f(1)=1$. By IVT, there exists $x^*$ with $f(x^*) = 0.6$. Then $v_A([0, x^*]) = 0.6$ and $v_B([x^*, 1]) = 1 - v_B([0, x^*])$. If $v_B([x^*,1]) \ge 0.4$, we are done. Otherwise, $v_B([0, x^*]) > 0.6$, so assign Bob $[0, x^*]$ (value $> 0.6 \ge 0.4$). By IVT, find $y^*$ with $v_B([0, y^*]) = 0.4$; then give Alice $[y^*, 1]$ with $v_A([y^*, 1]) = 1 - v_A([0, y^*])$. Since $y^* < x^*$ (Bob reaches $0.4$ before Alice reaches $0.6$), $v_A([y^*,1]) > 0.4$. The flexibility of choosing \emph{which} agent gets which side guarantees existence.
\end{exbox}

\separator

\begin{exbox}[Example 3: Price of Proportionality]
\textbf{Problem (Exercise 3).} $n = k^2$ agents. Agents $1, \ldots, k$ have valuation uniform on $[(i-1)/k,\, i/k]$ and $0$ elsewhere (agent $i$ values subintervals of $[(i-1)/k, i/k]$ at $k$ times their length). The remaining $k^2 - k$ agents value the entire cake uniformly.

\textbf{(a) Maximum social welfare.}
Give each ``specialist'' agent $i$ (for $i = 1, \ldots, k$) their favourite interval $[(i-1)/k, i/k]$. Each gets value $1$. The remaining $k^2 - k$ uniform agents get nothing (value $0$).
\[
  \text{Max SW} = k \cdot 1 + (k^2 - k) \cdot 0 = k.
\]

\textbf{(b) Maximum social welfare under proportionality.}
Each of the $n = k^2$ agents must get value $\ge 1/k^2$. The $k^2 - k$ uniform agents need total cake length $\ge (k^2 - k)/k^2 = 1 - 1/k$. This leaves $\le 1/k$ of the cake for the $k$ specialist agents combined. Since the specialists' favourite intervals are disjoint and each has length $1/k$, at most one specialist can get their full interval (impossible---we've only $1/k$ total). In fact, each specialist must get at least $1/k^2$ (proportional share) from their interval, contributing $k \cdot (1/k^2) = 1/k$ to specialist welfare at minimum.

The uniform agents' total welfare from their share: each gets $\ge 1/k^2$ cake length, so total uniform welfare $\ge (k^2 - k)/k^2$. The specialists' welfare: each specialist $i$ gets some portion of $[(i-1)/k, i/k]$ with total specialist cake $\le 1/k$, so total specialist welfare $\le k \cdot 1 = k$ is not achievable.

Optimally: give each specialist exactly $1/k^2$ of their interval (satisfying proportionality, value $= k \cdot 1/k^2 = 1/k$), and distribute the rest ($1 - k \cdot 1/k^2 = 1 - 1/k$) to uniform agents. Uniform agents sharing $1 - 1/k$ cake, each needing $\ge 1/k^2$, total uniform welfare $= 1 - 1/k$.

\[
  \text{Max proportional SW} = k \cdot \frac{1}{k} + (1 - \frac{1}{k}) = 1 + 1 - \frac{1}{k} = 2 - \frac{1}{k}.
\]

\textbf{Price of proportionality} $= k / (2 - 1/k) = k^2 / (2k - 1) \approx k/2$ for large $k$, i.e.\ $\approx \sqrt{n}/2$.
\end{exbox}

\separator

\begin{exbox}[Example 4: Cut-and-Choose Trace]
\textbf{Problem.} Cake: $[0, 1]$, left half strawberry, right half vanilla. Agent~A values strawberry at $0.8$, vanilla at $0.2$. Agent~B values both equally. Agent~A cuts. Where do they cut, and what does each agent get?

\textbf{Solution.} Agent~A wants two pieces of equal value ($0.5$ each). Strawberry density for A: $0.8/0.5 = 1.6$ per unit length on $[0, 0.5]$; vanilla density: $0.2/0.5 = 0.4$ on $[0.5, 1]$.

Cut at point $x$. Value of $[0, x]$ to A:
\begin{itemize}
  \item If $x \le 0.5$: $v_A = 1.6x$. Set $= 0.5$: $x = 5/16 = 0.3125$.
\end{itemize}
So A cuts at $x = 0.3125$. Piece~1: $[0, 0.3125]$, Piece~2: $[0.3125, 1]$.

Agent~B (uniform): $v_B(\text{Piece 1}) = 0.3125$, $v_B(\text{Piece 2}) = 0.6875$. B picks Piece~2.

\textbf{Result:} A gets $[0, 0.3125]$ (value $0.5$), B gets $[0.3125, 1]$ (value $0.6875$). Envy-free. Not equitable.
\end{exbox}

\separator

\begin{exbox}[Example 5: Even--Paz Trace]
\textbf{Problem.} Four agents with the following valuations on $[0,1]$: Agent~1 uniform, Agent~2 values $[0, 0.5]$ at $1$, Agent~3 values $[0.5, 1]$ at $1$, Agent~4 uniform. Run Even--Paz.

\textbf{Solution.} Each agent marks the midpoint of their value on $[0, 1]$:
\begin{itemize}
  \item Agent~1: $z_1 = 0.5$ (uniform $\Rightarrow$ midpoint by length).
  \item Agent~2: $z_2 = 0.25$ ($v_2([0, 0.25]) = 0.5$).
  \item Agent~3: $z_3 = 0.75$ ($v_3([0.5, 0.75]) = 0.5$, but need $v_3([0, z_3]) = 0.5$; since $v_3([0, 0.5]) = 0$, need $v_3([0.5, z_3]) = 0.5 \Rightarrow z_3 = 0.75$).
  \item Agent~4: $z_4 = 0.5$.
\end{itemize}
Sorted marks: $0.25, 0.5, 0.5, 0.75$. Median split $z^* = 0.5$. Left group (marks $\le 0.5$): Agents~$\{2, 1, 4\}$---but we need exactly 2 per side. Take the two smallest marks: Agents~2 and 1 go left on $[0, 0.5]$; Agents~4 and 3 go right on $[0.5, 1]$.

Recurse on each half with 2 agents each (essentially Cut-and-Choose). Each agent ends up with $\ge 1/4$ of total value.
\end{exbox}

\separator

\begin{exbox}[Example 6: Selfridge--Conway Trace]
\textbf{Problem.} Three agents. Agent~1 values the cake uniformly. Agent~2 values $[0, 0.5]$ at $0.8$ and $[0.5, 1]$ at $0.2$. Agent~3 values $[0, 0.5]$ at $0.3$ and $[0.5, 1]$ at $0.7$.

\textbf{Stage 0:} Agent~1 cuts into three equal pieces: $[0, 1/3], [1/3, 2/3], [2/3, 1]$, each worth $1/3$ to Agent~1.

Agent~2's valuations: $v_2([0, 1/3]) = 0.8 \cdot (2/3) = 0.533$, $v_2([1/3, 1/2]) = 0.8 \cdot (1/3) = 0.267$, so $v_2([1/3, 2/3]) = 0.267 + 0.2 \cdot (1/6) = 0.267 + 0.033 = 0.3$, $v_2([2/3, 1]) = 0.2 \cdot (1/3)/0.5 = 0.133$. Largest piece for Agent~2: $[0, 1/3]$ at $0.533$, second largest: $[1/3, 2/3]$ at $0.3$.

Agent~2 trims $[0, 1/3]$ to match $0.3$. The trimming (Cake~2) has value $0.533 - 0.3 = 0.233$ to Agent~2.

\textbf{Stage 1:} Agent~3 picks their favourite from the three Cake~1 pieces. Then allocation proceeds per the algorithm. \textbf{Stage 2:} Cake~2 divided by $T'$, chosen in order $T, 1, T'$. Result is envy-free.
\end{exbox}
